% 5_discussion.tex
\section{Discussion}
\label{sec:discussion}

\paragraph{A case-study-derived reporting heuristic.}
Our results support a concrete minimum for papers that derive
component-importance claims from LLM-judged multi-agent artifacts:
(1)~score with at least two judges from different model families, and
preferably different vendors; (2)~report pairwise judge agreement and
per-judge offsets alongside the scores; (3)~when judges disagree on the
ordering of the variants, downgrade the component claim to
``undetermined'' or escalate to human judgment, rather than reporting
the single judge's $p$-value. Condition (3) is the escalation idea
of~\cite{jung2024trust} applied at the level of conclusions, and (1)--(2)
are the panel idea of~\cite{verga2024juries} reduced to its
measurement core. We offer this as a heuristic derived from one case study: what the
study establishes is the need to report judge sensitivity, not that
this particular panel size is sufficient.

\paragraph{What the leniency result changes.}
The community's default mitigation for judge bias is family
substitution. Our decomposition shows why that is insufficient here:
the operative effect is a between-judge level shift (offset interval
excluding zero, persisting on transcripts the judge did not produce)
that survives family substitution, while self-preference---the effect
family substitution targets---is not detected at the available power. A
judge panel with a level-calibration step (e.g., anchoring all judges
to a common reference set) is the natural repair; we leave its
validation to future work.

\paragraph{Statistical honesty at this sample size.}
With five repetitions per ablation condition (three in the same-machine
replication), no contrast survives multiplicity correction, and the
$3\times3$ exact test cannot return $p<0.1$ at all. We therefore make
no significance claim for any single contrast. The evidence for RQ3 is
structural: on identical transcripts, changing only the judge changes
the direction of the component effect, with effect sizes up to
$\lvert\delta\rvert=1.00$; and the audited evaluation's reported
magnitudes ($-0.92$, $-0.80$) are not approached by any judge we
tested. Direction reversal under a controlled factor is a weaker but
more robust form of evidence than a $p$-value at this $n$, and it is
sufficient to break the specific inferential move under audit: reading
a single judge's component differences as properties of the system.

\paragraph{Relationship to general judge-reliability findings.}
That LLM judges are imperfect is established~\cite{reliability2026,
allverdicts2026}. The contribution here is downstream: measuring what
that imperfection does to a specific scientific inference (component
ranking), in a setting where no ground truth anchors the scores, and
separating two mechanisms that prior work treats jointly. The
same-machine replication holds transcripts fixed so that judge identity
is the only changed factor within the tested set; because the
factorial's driver and judge factors are partially confounded by
design, we do not claim a driver-independent causal effect beyond the
tested transcripts.

\paragraph{A defect found in the audited judge.}
The audit surfaced a concrete defect in the audited evaluation's judge
prompt: the pre-computed speaking-count statistic omits agents that
never speak, understating the maximum gap and misapplying the prompt's
own hard caps in 7 of 32 discussions. We preserved the audited behavior
in the main analysis (the prompt is part of the protocol under audit,
and the defect applies identically to every judge, so it cannot
manufacture inter-judge disagreement) and verified by re-judgment that
the affected discussions do not carry the qualitative conclusions
(Appendix~\ref{app:robust}). Reporting the defect is part of the audit's
value: it is exactly the kind of silent evaluation bug that
single-judge protocols do not surface.
